\clearpage\pagestyle{otMain}\setlength{\otBlocEspacement}{2.5mm}\fontsize{11}{15}\selectfont\corrigesChapitre{7}{Comparer et additionner des fractions}
\begin{blocLivre}{cor-7-1}\begin{corrige}{1}{Mêmes parts}\demarche{À parts égales, comparer les numérateurs.}
$\res{\frac79>\frac59}$ : mêmes parts, plus de parts. Ordre : $\res{\frac2{15}<\frac8{15}<\frac{11}{15}}$.
\end{corrige}\end{blocLivre}
\begin{blocLivre}{cor-7-2}\begin{corrige}{2}{Mêmes numérateurs}\demarche{À numérateur égal, le plus petit dénominateur l’emporte.}
Trois parts chacune, mais les cinquièmes sont plus grands : $\res{\frac35>\frac38}$. De même $\res{\frac7{10}>\frac7{12}}$.
\end{corrige}\end{blocLivre}
\begin{blocLivre}{cor-7-3}\begin{corrige}{3}{Écritures équivalentes}\demarche{Multiplier haut et bas par le même nombre.}
Les parts deviennent plus fines : $\frac56=\res{\frac{10}{12}}$ ; $\frac34=\res{\frac{15}{20}}$ ; $\frac23=\res{\frac{40}{60}}$.
\end{corrige}\end{blocLivre}
\begin{blocLivre}{cor-7-4}\begin{corrige}{4}{Comparer en convertissant}\demarche{Convertir vers le dénominateur multiple.}
$\frac56=\frac{10}{12}<\frac{11}{12}$ : $\res{\frac56<\frac{11}{12}}$. $\frac23=\frac8{12}>\frac7{12}$ : $\res{\frac23>\frac7{12}}$.
\end{corrige}\end{blocLivre}
\begin{blocLivre}{cor-7-5}\begin{corrige}{5}{Autour de un demi}\demarche{Écrire un demi avec le dénominateur utile.}
$\frac12=\frac48$ : $\res{\frac38<\frac12}$. $\frac12=\frac6{12}$ : $\res{\frac7{12}>\frac12}$. $\frac5{10}=\res{\frac12}$ exactement.
\end{corrige}\end{blocLivre}
\begin{blocLivre}{cor-7-6}\begin{corrige}{6}{Sommes directes}\demarche{Garder le dénominateur commun.}
$\res{\frac9{12}}$ et $\res{\frac7{15}}$ : les parts s’ajoutent ou se retirent, le partage reste.
\end{corrige}\end{blocLivre}
\begin{blocLivre}{cor-7-7}\begin{corrige}{7}{Convertir puis additionner}\demarche{Convertir vers le dénominateur multiple.}
$\frac34=\frac68$ : $\frac68+\frac58=\res{\frac{11}8}=1+\frac38$. $\frac23=\frac46$ : $\frac46+\frac16=\res{\frac56}$.
\end{corrige}\end{blocLivre}
\begin{blocLivre}{cor-7-8}\begin{corrige}{8}{Retirer à un entier}\demarche{Écrire l’entier avec le dénominateur utile.}
$1=\frac{15}{15}$ : il reste $\res{\frac{11}{15}}$. $2=\frac{12}6$ : il reste $\frac{12}6-\frac56=\res{\frac76}$, soit $1+\frac16$.
\end{corrige}\end{blocLivre}
\begin{blocLivre}{cor-7-9}\begin{corrige}{9}{Compléter}\demarche{Penser « ce qui manque » ou « défaire ».}
$\square=\res{\frac3{10}}$ ; $\square=\res{\frac48}$ ; $\square=\frac1{12}+\frac5{12}=\res{\frac6{12}}$, soit un demi.
\end{corrige}\end{blocLivre}
\begin{blocLivre}{cor-7-10}\begin{corrige}{10}{L’erreur classique}\demarche{Revenir aux parts d’une même unité.}
Réunir deux quarts donne deux parts d’un partage en 4, pas un partage en 8 : $\frac14+\frac14=\res{\frac24}$, la moitié. On n’additionne jamais les dénominateurs.
\end{corrige}\end{blocLivre}
\begin{blocLivre}{cor-7-11}\begin{corrige}{11}{La tarte partagée}\demarche{Convertir les quarts en huitièmes.}
$\frac14=\frac28$ : ensemble $\frac38+\frac28=\res{\frac58}$. Il reste $1-\frac58=\res{\frac38}$ de tarte. Contrôle : $\frac58+\frac38=1$.
\end{corrige}\end{blocLivre}
\begin{blocLivre}{cor-7-12}\begin{corrige}{12}{Le trajet en trois temps}\demarche{Réunir les parts connues avant le reste.}
$\frac14=\frac3{12}$ : bus et tram couvrent $\frac7{12}+\frac3{12}=\frac{10}{12}$. À pied : $1-\frac{10}{12}=\res{\frac2{12}}$, soit un sixième du trajet.
\end{corrige}\end{blocLivre}
\begin{blocLivre}{cor-7-13}\begin{corrige}{13}{Les bouteilles d’eau}\demarche{Comparer la somme à l’unité.}
$\frac34=\frac68$ : total $\frac68+\frac58=\frac{11}8=1+\frac38$. Il a bu $\res{\text{plus}}$ d’un litre, de $\res{\frac38}$ de litre.
\end{corrige}\end{blocLivre}
\begin{blocLivre}{cor-7-14}\begin{corrige}{14}{Le classement des lectures}\demarche{Choisir un dénominateur commun aux trois.}
En vingtièmes : $\frac7{10}=\frac{14}{20}$, $\frac35=\frac{12}{20}$, $\frac34=\frac{15}{20}$. Ordre croissant : $\res{\frac35<\frac7{10}<\frac34}$ : c’est $\res{\text{Chloé}}$ la plus avancée.
\end{corrige}\end{blocLivre}
\begin{blocLivre}{cor-7-15}\begin{corrige}{15}{Un pas vers l’abstrait}\demarche{Observer quand les numérateurs recomposent le dénominateur.}
Cinq parts et huit parts font treize parts d’un partage en treize : l’unité entière, $\res{\frac5{13}+\frac8{13}=1}$. De même $37+23=60$ : $\res{\frac{37}{60}+\frac{23}{60}=1}$.
\end{corrige}\end{blocLivre}
